% =====================================================================
%  c04.tex —— §5 扭积（twisted product）
%  原书 p.50–53；PDF p.9（§5 标题）–p.12（§6「反例」之前）
%  本块无编号公式（原件 (1)–(10) 中无一条落在本块）。
%  定理类环境：定义 5.1、例 5.2、例 5.3、定理 5.4
%  跨块引用一律写原书字面节号、定理号与文献号 [N]，不写 \ref / \cite。
% =====================================================================

% ===== 原书 p.50 =====
\section{\texorpdfstring{$\ainf$范畴的扭积}{A∞ 范畴的扭积}}
\label{sec:c04-5}

\begin{definition}\label{def:c04-5.1}
设$C$为某个$\ainf$范畴。设给定$C$的两个单侧扭复形
$(\mathcal{X}^{-},\Phi^{-})$、$(\mathcal{X}^{+},\Phi^{+})$，满足
\[
\mathcal{X}^{-}=\mathcal{X}\oplus\mathcal{Y},\qquad
\mathcal{X}^{+}=\mathcal{Y}\oplus\mathcal{Z},
\]
\[
\Phi^{-}=
\begin{pmatrix}
\Phi_{\mathcal{X}} & \Phi_{\mathcal{X}\mathcal{Y}}\\
0 & \Phi_{\mathcal{Y}}
\end{pmatrix},\qquad
\Phi^{+}=
\begin{pmatrix}
\Phi_{\mathcal{Y}} & \Phi_{\mathcal{Y}\mathcal{Z}}\\
0 & \Phi_{\mathcal{Z}}
\end{pmatrix}.
\]
其中$\mathcal{X},\mathcal{Y},\mathcal{Z}$都不是零对象。令
\[
\Phi:=
\begin{pmatrix}
\Phi_{\mathcal{X}} & \Phi_{\mathcal{X}\mathcal{Y}} & 0\\
0 & \Phi_{\mathcal{Y}} & \Phi_{\mathcal{Y}\mathcal{Z}}\\
0 & 0 & \Phi_{\mathcal{Z}}
\end{pmatrix},
\]
并把
\[
\tilde{m}_{*}(\Phi):=\tilde{m}_{1}(\Phi)+\tilde{m}_{2}(\Phi,\Phi)
+\tilde{m}_{3}(\Phi,\Phi,\Phi)+\cdots
\]
重新按$\mathcal{X}\oplus\mathcal{Y}\oplus\mathcal{Z}$分块分解。所得$\mathcal{X}$-$\mathcal{Z}$块分量
（右上角的块分量）称为$C$中的
\textbf{扭积}$(\mathcal{X},\mathcal{Y},\mathcal{Z};\Phi)$。
\end{definition}

\transnote{定义中的$\mathcal{X},\mathcal{Y},\mathcal{Z}$是$\tilde{C}$的对象，即$C$中
对象的平移的形式直和，一般为若干对象之直和，而不是$C$的单个对象。「不是零对象」
只为排除退化情形。两个单侧扭复形通过公共部分$\mathcal{Y}$拼成一个$3\times3$的
单侧扭复形$\Phi$，扭积取它在$(\mathcal{X},\mathcal{Z})$位置的分量。}

由于$(\mathcal{X}^{-},\Phi^{-})$、$(\mathcal{X}^{+},\Phi^{+})$是单侧扭复形，
$\Phi_{\mathcal{X}}$、$\Phi_{\mathcal{Y}}$、$\Phi_{\mathcal{Z}}$都是对角元为零的上三角
矩阵，$\Phi$也是如此。特别地，$(\mathcal{X},\Phi_{\mathcal{X}})$、
$(\mathcal{Y},\Phi_{\mathcal{Y}})$、$(\mathcal{Z},\Phi_{\mathcal{Z}})$是单侧扭
复形，$\Phi_{\mathcal{X}\mathcal{Y}}$是从$(\mathcal{X},\Phi_{\mathcal{X}})$到
$(\mathcal{Y},\Phi_{\mathcal{Y}})$的$m_{1}^{Tw}$-闭形式，
$\Phi_{\mathcal{Y}\mathcal{Z}}$是从$(\mathcal{Y},\Phi_{\mathcal{Y}})$到
$(\mathcal{Z},\Phi_{\mathcal{Z}})$的$m_{1}^{Tw}$-闭形式。因此，扭积
$(\mathcal{X},\mathcal{Y},\mathcal{Z};\Phi)$无非是$s\Tw(\mathcal{C})$中的积
$m_{2}^{Tw}(\Phi_{\mathcal{X}\mathcal{Y}},\Phi_{\mathcal{Y}\mathcal{Z}})$。于是
$\tilde{m}_{*}(\Phi)$在$\mathcal{X}$-$\mathcal{Z}$块分量以外均为平凡，扭积
$(\mathcal{X},\mathcal{Y},\mathcal{Z};\Phi)$是从$(\mathcal{X},\Phi_{\mathcal{X}})$
到$(\mathcal{Z},\Phi_{\mathcal{Z}})$的$m_{1}^{Tw}$-闭形式。

这些扭积包含 DG 范畴的 Massey 积，以及它到$\ainf$范畴情形的自然推广。下面用简单
的例子来看。

\begin{example}[{DG 范畴的 Massey 积（cf.\ Gelfand-Manin [5]）}]
\label{ex:c04-5.2}
设$C$为 DG 范畴。对对象$X_{1},X_{2},X_{3},X_{4}\in C$，取
$\tilde{X}_{1},\tilde{X}_{2},\tilde{X}_{3},\tilde{X}_{4}\in\tilde{C}$为把它们的
次数作若干平移所得的对象，即设$\tilde{X}_{i}:=X_{i}[r_{i}]$，$r_{i}\in\Z$，
$i=1,2,3,4$。考虑零次的$\tilde{m}_{1}$-闭形式
\[
\rho_{i(i+1)}\in(s\tilde{C})^{0}(\tilde{X}_{i},\tilde{X}_{i+1}),
\qquad i=1,2,3,
\]
使它们在上同调上满足
\[
[\rho_{12}\cdot\rho_{23}]=0,\qquad [\rho_{23}\cdot\rho_{34}]=0.
\]
其中$\cdot$表示 DG 范畴$\tilde{C}$中的积。取态射
$h_{13}\in s\tilde{C}^{0}(\tilde{X}_{1},\tilde{X}_{3})$、
$h_{24}\in s\tilde{C}^{0}(\tilde{X}_{2},\tilde{X}_{4})$，使得
\[
d(h_{13})=-\rho_{12}\cdot\rho_{23},\qquad
d(h_{24})=-\rho_{23}\cdot\rho_{34}=0.
\]

% ===== 原书 p.51 =====
此时，考虑$C$的扭积
\[
(\tilde{X}_{1},\tilde{X}_{2}\oplus\tilde{X}_{3},\tilde{X}_{4};\Phi),
\qquad
\Phi:=
\begin{pmatrix}
0 & \rho_{12} & h_{13} & 0\\
0 & 0 & \rho_{23} & h_{24}\\
0 & 0 & 0 & \rho_{34}\\
0 & 0 & 0 & 0
\end{pmatrix}.
\]
由
\[
\begin{aligned}
\tilde{m}_{*}(\Phi)&=d(\Phi)+\Phi\cdot\Phi\\
&=\begin{pmatrix}
0&0&d(h_{13})&0\\
0&0&0&d(h_{24})\\
0&0&0&0\\
0&0&0&0
\end{pmatrix}
+\begin{pmatrix}
0&0&\rho_{12}\cdot\rho_{23}&h_{13}\cdot\rho_{34}+\rho_{12}\cdot h_{24}\\
0&0&0&\rho_{23}\cdot\rho_{34}\\
0&0&0&0\\
0&0&0&0
\end{pmatrix}\\
&=\begin{pmatrix}
0&0&0&h_{13}\cdot\rho_{34}+\rho_{12}\cdot h_{24}\\
0&0&0&0\\
0&0&0&0\\
0&0&0&0
\end{pmatrix},
\end{aligned}
\]
可知，作为$(1,4)$分量得到的扭积无非是 Massey 积。
\end{example}

\transnote{$d(h_{13})=-\rho_{12}\cdot\rho_{23}$的负号不是笔误：它来自
$\tilde{m}_{*}(\Phi)=d(\Phi)+\Phi\cdot\Phi=0$的$(1,3)$分量，$(2,4)$分量同理给出
$d(h_{24})=-\rho_{23}\cdot\rho_{34}$。原文在后一式又添了「$=0$」，与上文「在上同调上
$[\rho_{23}\cdot\rho_{34}]=0$」的措辞不一致：按字面二者不能同时成立。}

下面看上一例在$\ainf$范畴情形的对应物。

\begin{example}[{（极小）$\ainf$范畴中的 Massey 积}]
\label{ex:c04-5.3}
设$C$为极小$\ainf$范畴。此时注意：态射空间的上同调与态射空间自身一致。与例 5.2
相同，对对象$X_{1},X_{2},X_{3},X_{4}\in\tilde{C}$，取
$\tilde{X}_{1},\tilde{X}_{2},\tilde{X}_{3},\tilde{X}_{4}\in\tilde{C}$为把它们的
次数作若干平移所得的对象。设态射
$\rho_{i(i+1)}\in s\tilde{C}^{0}(\tilde{X}_{i},\tilde{X}_{i+1})$，$i=1,2,3$，满足
\[
\rho_{12}\cdot\rho_{23}=0,\qquad \rho_{23}\cdot\rho_{34}=0.
\]
此时，考虑扭积
\[
(\tilde{X}_{1},\tilde{X}_{2}\oplus\tilde{X}_{3},\tilde{X}_{4};\Phi),
\qquad
\Phi=
\begin{pmatrix}
0 & \rho_{12} & 0 & 0\\
0 & 0 & \rho_{23} & 0\\
0 & 0 & 0 & \rho_{34}\\
0 & 0 & 0 & 0
\end{pmatrix}.
\]
由
\[
\tilde{m}_{*}(\Phi)=\tilde{m}_{3}(\Phi,\Phi,\Phi)=
\begin{pmatrix}
0&0&0&\tilde{m}_{3}(\rho_{12},\rho_{23},\rho_{34})\\
0&0&0&0\\
0&0&0&0\\
0&0&0&0
\end{pmatrix},
\]

% ===== 原书 p.52 =====
得到扭积$\tilde{m}_{3}(\rho_{12},\rho_{23},\rho_{34})$。
\end{example}

\transnote{极小范畴中$m_{1}=0$，故$\tilde{m}_{1}=0$。与例 5.2 相比，这里
$\rho_{12}\cdot\rho_{23}=0$、$\rho_{23}\cdot\rho_{34}=0$都是严格等式，不需要
$h_{13},h_{24}$这样的提升，于是$\tilde{m}_{2}(\Phi,\Phi)=0$。又$\Phi$的非零元只有
三条次对角元，$n\ge4$的$\tilde{m}_{n}(\Phi,\dots,\Phi)$必含零元，故扭积只剩
$\tilde{m}_{3}(\Phi,\Phi,\Phi)$一项。这正是极小范畴与一般 DG 范畴的差别。}

把这看作 Massey 积在$\ainf$范畴情形的对应物是自然的。事实上，设例 5.3 中的极小
$\ainf$范畴$C$为例 5.2 中 DG 范畴$C$的极小模型。特别地，如定理 2.3 所述，极小
模型有具体的构造方法。若例 5.3 中的极小$\ainf$范畴$C$是由该方法从例 5.2 中的 DG
范畴$C$构造出的极小$\ainf$范畴，则在这一等同之下，例 5.2、5.3 中的扭积一致
[12]。

例 5.2、5.3 两种情形下，扭积
$(\tilde{X}_{1},\tilde{X}_{2}\oplus\tilde{X}_{3},\tilde{X}_{4};\Phi)$对应的都是
$\Tr(\mathcal{C})$中的态射合成
\[
[\rho_{12}\cdot\rho_{34}]:(\tilde{X}_{1},0)\longrightarrow
(\tilde{X}_{23},\Phi_{\tilde{X}_{23}})\longrightarrow(\tilde{X}_{4},0),
\]
其中
\[
\tilde{X}_{23}:=\tilde{X}_{2}\oplus\tilde{X}_{3},\qquad
\Phi_{\tilde{X}_{23}}:=
\begin{pmatrix}0&\rho_{12}\\0&0\end{pmatrix}.
\]

\transnote{扭积落在
$(s\Tw(\mathcal{C}))^{1}((\mathcal{X},\Phi_{\mathcal{X}}),(\mathcal{Z},\Phi_{\mathcal{Z}}))$
中，取$m_{1}^{Tw}$-上同调类即得$\Tr(\mathcal{C})$的态射，所以它对应
$\Tr(\mathcal{C})$里的一段合成：两段态射分别由$\rho_{12}$、$\rho_{34}$给出，中间扭复形
$(\tilde{X}_{23},\Phi_{\tilde{X}_{23}})$的扭部分记录必要的中间项。原文把
$\Phi_{\tilde{X}_{23}}$的$(1,2)$分量印作$\rho_{12}$，但它在例 5.2 的$4\times4$矩阵中
对应$(2,3)$分量$\rho_{23}$；译文照录原文。}

扭积还包含 Massey 积的推广——高阶 Massey 积，以及它的进一步推广——矩阵
Massey 积（[17]）；这些包含关系见 [10, 12]。有些扭积甚至不对应矩阵 Massey 积，但即便如此，
也不能从扭积完全还原原来的$\ainf$结构。后文构造的问题 4.1 的反例（例 6.1）中，
将构造这样的高阶$\ainf$积：它对任何扭积都不（按下述意义）显式地有贡献，却无法
通过$\ainf$等价消去。

扭积$(\mathcal{X},\mathcal{Y},\mathcal{Z};\Phi)$是
$(s\Tw(\mathcal{C}))^{1}((\mathcal{X},\Phi_{\mathcal{X}}),
(\mathcal{Z},\Phi_{\mathcal{Z}}))$中的$m_{1}^{Tw}$-闭形式。为从这些闭形式读出
$\Tr(\mathcal{C})$的乘积结构，把
$(s\Tw(\mathcal{C}))^{1}((\mathcal{X},\Phi_{\mathcal{X}}),
(\mathcal{Z},\Phi_{\mathcal{Z}}))$中$m_{1}^{Tw}$-恰当形式构成的空间记为
\[
B^{1}((\mathcal{X},\Phi_{\mathcal{X}}),(\mathcal{Z},\Phi_{\mathcal{Z}}))=:B^{1}.
\]

由扭积的定义得到以下结论。

\begin{theorem}\label{thm:c04-5.4}
设$C$、$C'$为具有相同对象与态射的$\ainf$范畴（即
$\Ob(C)=\Ob(C')$，且对任意$X,Y\in\Ob(C)=\Ob(C')$有
$C^{i}(X,Y)\simeq C^{\prime i}(X,Y)$）。若$C$、$C'$满足以下条件，则$C$与$C'$的
对象和态射的等同给出范畴同构$C\simeq C'$，并且这一等同还扩展为三角范畴同构
$\Tr(\mathcal{C})\simeq\Tr(\mathcal{C}')$。
\begin{itemize}[leftmargin=1.6em,itemsep=2pt,topsep=3pt]
\item $(\mathcal{X},\Phi_{\mathcal{X}})$在$C$中是单侧扭复形，当且仅当它在$C'$中
是单侧扭复形。
\item 对任意两个单侧扭复形$(\mathcal{X},\Phi_{\mathcal{X}})$、
$(\mathcal{Z},\Phi_{\mathcal{Z}})$，$B^{1}$在$C$中的情形与在$C'$中的情形一致。
\item 对任意扭积$(\mathcal{X},\mathcal{Y},\mathcal{Z};\Phi)$，它在$C$中的情形与
在$C'$中的情形除$B^{1}$之外一致。
\end{itemize}
\hfill$\square$
\end{theorem}

% ===== 原书 p.53 =====
因此，只要构造出两个满足这一定理条件的极小$\ainf$范畴$C$、$C'$，并且它们互不
$\ainf$同型，那么由定理 4.2，它们就是问题 4.1 的反例。
